% Authored: build interface and packaging.

\chapter{Build and Packaging}
\label{ch:build}

\section{Build interface}

The Makefile is the documented interface. Every target is a thin wrapper over a
Poetry command, so nothing is hidden behind Make-specific logic and every
target can be run directly if Make is unavailable.

% !TEX root = ../main.tex
% ---------------------------------------------------------------
% GENERATED FILE - DO NOT EDIT.
% Produced by docs/tools/render_engineering.py from the repository source.
% Regenerate with: make docs
% ---------------------------------------------------------------

\begin{longtable}{>{\raggedright\arraybackslash\small}p{0.18\textwidth} >{\raggedright\arraybackslash\small}p{0.24\textwidth} >{\raggedright\arraybackslash\small}p{0.50\textwidth}}
\caption{Makefile targets: the documented build interface.}\label{tab:make-targets}\\
\toprule
Target & Prerequisites & Recipe \\
\midrule
\endfirsthead
\toprule
Target & Prerequisites & Recipe \\
\midrule
\endhead
\midrule
\multicolumn{3}{r}{\footnotesize continued on next page}\\
\endfoot
\bottomrule
\endlastfoot
\texttt{install} & -- & \texttt{poetry install} \\
\texttt{test} & -- & \texttt{poetry run pytest} \\
\texttt{lint} & -- & \texttt{poetry run ruff check .\allowbreak{}} \\
\texttt{format-\allowbreak{}check} & -- & \texttt{poetry run ruff format -\allowbreak{}-\allowbreak{}check .\allowbreak{}} \\
\texttt{typecheck} & -- & \texttt{poetry run mypy src} \\
\texttt{typecheck-\allowbreak{}docs} & -- & \texttt{poetry run mypy -\allowbreak{}-\allowbreak{}python-\allowbreak{}version \$(DOCS\_\allowbreak{}PYTHON\_\allowbreak{}VERSION)\allowbreak{} docs/\allowbreak{}tools} \\
\texttt{coverage} & -- & \texttt{poetry run pytest -\allowbreak{}-\allowbreak{}cov=feedback\_\allowbreak{}intelligence\_\allowbreak{}agent -\allowbreak{}-\allowbreak{}cov-\allowbreak{}report=term-\allowbreak{}missing -\allowbreak{}-\allowbreak{}cov-\allowbreak{}fail-\allowbreak{}under=\$(COVERAGE\_\allowbreak{}THRESHOLD)\allowbreak{}} \\
\texttt{build} & -- & \texttt{poetry build} \\
\texttt{quality} & \texttt{lint}, \texttt{typecheck}, \texttt{typecheck-\allowbreak{}docs}, \texttt{test} & -- \\
\texttt{ci} & \texttt{lint}, \texttt{format-\allowbreak{}check}, \texttt{typecheck}, \texttt{typecheck-\allowbreak{}docs}, \texttt{coverage}, \texttt{build} & -- \\
\texttt{docs} & -- & \texttt{\$(DOCS\_\allowbreak{}ENV)\allowbreak{} poetry run python -\allowbreak{}m tools.\allowbreak{}generate\_\allowbreak{}docs -\allowbreak{}-\allowbreak{}repo-\allowbreak{}root .\allowbreak{} \textbackslash{} \&\& \$(if \$(REVISION)\allowbreak{},\allowbreak{}-\allowbreak{}-\allowbreak{}revision \$(REVISION)\allowbreak{},\allowbreak{})\allowbreak{}} \\
\texttt{docs-\allowbreak{}check} & -- & \texttt{\$(DOCS\_\allowbreak{}ENV)\allowbreak{} poetry run python -\allowbreak{}m tools.\allowbreak{}validate\_\allowbreak{}docs -\allowbreak{}-\allowbreak{}repo-\allowbreak{}root .\allowbreak{}} \\
\texttt{docs-\allowbreak{}pdf} & \texttt{docs} & \texttt{cd \$(LATEX\_\allowbreak{}DIR)\allowbreak{} \&\& \$(LATEXMK)\allowbreak{} main.\allowbreak{}tex \textbar{}\textbar{} true \&\& cd \$(LATEX\_\allowbreak{}DIR)\allowbreak{} \&\& \$(LATEXMK)\allowbreak{} main.\allowbreak{}tex \textbar{}\textbar{} true \&\& \$(DOCS\_\allowbreak{}ENV)\allowbreak{} poetry run python -\allowbreak{}m tools.\allowbreak{}check\_\allowbreak{}latex\_\allowbreak{}log \$(LATEX\_\allowbreak{}DIR)\allowbreak{}/\allowbreak{}main.\allowbreak{}log} \\
\texttt{docs-\allowbreak{}pdf-\allowbreak{}reference} & \texttt{docs} & \texttt{cd \$(LATEX\_\allowbreak{}DIR)\allowbreak{} \&\& \$(LATEXMK)\allowbreak{} reference.\allowbreak{}tex \textbar{}\textbar{} true \&\& cd \$(LATEX\_\allowbreak{}DIR)\allowbreak{} \&\& \$(LATEXMK)\allowbreak{} reference.\allowbreak{}tex \textbar{}\textbar{} true \&\& \$(DOCS\_\allowbreak{}ENV)\allowbreak{} poetry run python -\allowbreak{}m tools.\allowbreak{}check\_\allowbreak{}latex\_\allowbreak{}log \$(LATEX\_\allowbreak{}DIR)\allowbreak{}/\allowbreak{}reference.\allowbreak{}log} \\
\texttt{docs-\allowbreak{}clean} & -- & \texttt{cd \$(LATEX\_\allowbreak{}DIR)\allowbreak{} \&\& latexmk -\allowbreak{}C main.\allowbreak{}tex reference.\allowbreak{}tex} \\
\texttt{demo} & -- & \texttt{poetry run python scripts/\allowbreak{}run\_\allowbreak{}demo.\allowbreak{}py} \\
\texttt{api} & -- & \texttt{poetry run uvicorn feedback\_\allowbreak{}intelligence\_\allowbreak{}agent.\allowbreak{}api:create\_\allowbreak{}app -\allowbreak{}-\allowbreak{}factory -\allowbreak{}-\allowbreak{}reload} \\
\texttt{docker-\allowbreak{}build} & -- & \texttt{docker build -\allowbreak{}t feedback-\allowbreak{}intelligence-\allowbreak{}agent .\allowbreak{}} \\
\texttt{docker-\allowbreak{}run} & -- & \texttt{docker run -\allowbreak{}-\allowbreak{}rm -\allowbreak{}p 8000:8000 feedback-\allowbreak{}intelligence-\allowbreak{}agent} \\
\texttt{compose-\allowbreak{}prod-\allowbreak{}up} & -- & \texttt{docker compose -\allowbreak{}f deploy/\allowbreak{}docker-\allowbreak{}compose.\allowbreak{}prod.\allowbreak{}yml up -\allowbreak{}d} \\
\texttt{compose-\allowbreak{}prod-\allowbreak{}down} & -- & \texttt{docker compose -\allowbreak{}f deploy/\allowbreak{}docker-\allowbreak{}compose.\allowbreak{}prod.\allowbreak{}yml down} \\
\texttt{clean} & -- & \texttt{rm -\allowbreak{}rf .\allowbreak{}artifacts .\allowbreak{}pytest\_\allowbreak{}cache .\allowbreak{}mypy\_\allowbreak{}cache .\allowbreak{}ruff\_\allowbreak{}cache htmlcov .\allowbreak{}coverage dist build} \\
\end{longtable}


\section{Python packaging}

The distribution is built by \texttt{poetry-core} from the \texttt{src} layout:
\texttt{packages = [\{ include = "feedback\_intelligence\_agent", from = "src" \}]}.
\texttt{poetry build} produces both an sdist and a wheel into \srcfile{dist/}.
The build runs on every push as the final CI step, so packaging breakage
surfaces immediately rather than at release time.

One console script is declared:
\texttt{feedback-agent = "feedback\_intelligence\_agent.cli:app"}.

Version numbers appear in three places and must agree: the
\texttt{version} field in \srcfile{pyproject.toml}, the
\texttt{\_\_version\_\_} attribute in the package \texttt{\_\_init\_\_}, and
the \texttt{version} fields of \srcfile{CITATION.cff} and
\srcfile{.zenodo.json}. Nothing in the repository enforces the agreement; see
\cref{ch:improvements}.

\section{Container images}

% !TEX root = ../main.tex
% ---------------------------------------------------------------
% GENERATED FILE - DO NOT EDIT.
% Produced by docs/tools/render_engineering.py from the repository source.
% Regenerate with: make docs
% ---------------------------------------------------------------

\section{Container images}

\begin{longtable}{>{\raggedright\arraybackslash\small}p{0.22\textwidth} >{\raggedright\arraybackslash\small}p{0.24\textwidth} >{\raggedright\arraybackslash\small}p{0.10\textwidth} >{\raggedright\arraybackslash\small}p{0.36\textwidth}}
\caption{Container images defined in the repository.}\label{tab:container-images}\\
\toprule
Dockerfile & Base image(s) & Ports & Entrypoint \\
\midrule
\endfirsthead
\toprule
Dockerfile & Base image(s) & Ports & Entrypoint \\
\midrule
\endhead
\midrule
\multicolumn{4}{r}{\footnotesize continued on next page}\\
\endfoot
\bottomrule
\endlastfoot
\texttt{Dockerfile} & \texttt{python:3.\allowbreak{}11-\allowbreak{}slim} & 8000 & \texttt{["uvicorn",\allowbreak{} "feedback\_\allowbreak{}intelligence\_\allowbreak{}agent.\allowbreak{}api:create\_\allowbreak{}app",\allowbreak{} "-\allowbreak{}-\allowbreak{}factory",\allowbreak{} "-\allowbreak{}-\allowbreak{}host",\allowbreak{} "0.\allowbreak{}0.\allowbreak{}0.\allowbreak{}0",\allowbreak{} "-\allowbreak{}-\allowbreak{}port",\allowbreak{} "8000"]\allowbreak{}} \\
\texttt{frontend/\allowbreak{}Dockerfile} & \texttt{node:22-\allowbreak{}alpine}, \texttt{node:22-\allowbreak{}alpine} & 4173 & \texttt{["npm",\allowbreak{} "run",\allowbreak{} "preview"]\allowbreak{}} \\
\end{longtable}

\section{Compose stacks}

\begin{longtable}{>{\raggedright\arraybackslash\small}p{0.26\textwidth} >{\raggedright\arraybackslash\small}p{0.18\textwidth} >{\raggedright\arraybackslash\small}p{0.20\textwidth} >{\raggedright\arraybackslash\small}p{0.14\textwidth} >{\raggedright\arraybackslash\small}p{0.12\textwidth}}
\caption{Compose services.}\label{tab:compose-services}\\
\toprule
File & Service & Image & Ports & Healthcheck \\
\midrule
\endfirsthead
\toprule
File & Service & Image & Ports & Healthcheck \\
\midrule
\endhead
\midrule
\multicolumn{5}{r}{\footnotesize continued on next page}\\
\endfoot
\bottomrule
\endlastfoot
\texttt{docker-\allowbreak{}compose.\allowbreak{}yml} & \texttt{feedback-\allowbreak{}agent} & built locally & 8000:8000 & no \\
\texttt{docker-\allowbreak{}compose.\allowbreak{}yml} & \texttt{frontend} & built locally & 4173:4173 & no \\
\texttt{docker-\allowbreak{}compose.\allowbreak{}yml} & \texttt{qdrant} & \texttt{qdrant/\allowbreak{}qdrant:latest} & 6333:6333, 6334:6334 & no \\
\texttt{deploy/\allowbreak{}docker-\allowbreak{}compose.\allowbreak{}prod.\allowbreak{}yml} & \texttt{api} & \texttt{feedback-\allowbreak{}intelligence-\allowbreak{}agent:latest} & 8000:8000 & yes \\
\end{longtable}

\section{Deployment descriptors}

\begin{longtable}{>{\raggedright\arraybackslash\small}p{0.40\textwidth} >{\raggedright\arraybackslash\small}p{0.52\textwidth}}
\caption{Deployment descriptors present in the repository.}\label{tab:deployments}\\
\toprule
Descriptor & Target \\
\midrule
\endfirsthead
\toprule
Descriptor & Target \\
\midrule
\endhead
\midrule
\multicolumn{2}{r}{\footnotesize continued on next page}\\
\endfoot
\bottomrule
\endlastfoot
\texttt{deploy/\allowbreak{}docker-\allowbreak{}compose.\allowbreak{}prod.\allowbreak{}yml} & Production-like Docker Compose stack. \\
\texttt{deploy/\allowbreak{}ecs-\allowbreak{}fargate-\allowbreak{}task-\allowbreak{}definition.\allowbreak{}json} & AWS ECS Fargate task definition. \\
\texttt{deploy/\allowbreak{}fly.\allowbreak{}toml} & Fly.io application configuration. \\
\end{longtable}


The application image is a single-stage \texttt{python:3.11-slim} build that
installs Poetry at a pinned version, copies \srcfile{pyproject.toml},
\srcfile{poetry.lock}, \srcfile{README.md}, \srcfile{src}, \srcfile{data},
\srcfile{docs}, and \srcfile{scripts}, then installs the main dependency group
followed by the project itself. It sets
\texttt{POETRY\_VIRTUALENVS\_CREATE=false} so packages land in the system
environment, exposes port 8000, and runs Uvicorn against the application
factory.

Two properties are worth calling out because they are deliberate:

\begin{itemize}[nosep]
  \item \srcfile{data/} is copied into the image, so a container started with
    no configuration can build an index and answer questions offline --- the
    containerised expression of \cref{adr:001}.
  \item The lock file is copied before installation, so the image resolves to
    the locked dependency set rather than to whatever is current.
\end{itemize}

The frontend image is a two-stage \texttt{node:22-alpine} build that installs
with \texttt{npm ci}, builds, and serves the preview server on port 4173.

\section{Frontend build}

The browser client is built with Vite and TypeScript. Its scripts are declared
in \srcfile{frontend/package.json} and it is built in CI by the frontend
workflow and again by the release workflow, whose artifact upload makes the
built assets available for a release.
